\documentclass[sigconf,nonacm]{acmart}

\usepackage{booktabs}
\usepackage{amsmath}
\usepackage{amssymb}
\usepackage{amsfonts}
\usepackage{graphicx}
\usepackage{subcaption}
\usepackage{microtype}
\usepackage{hyperref}

\settopmatter{printacmref=false}
\renewcommand\footnotetextcopyrightpermission[1]{}
\pagestyle{plain}

\title[Resolving Felsenstein's Dilemma in Phenomics]{Resolving Felsenstein's Dilemma in Phenomics: Continuous Spectral Timetree Manifolds Establish the Predictability Frontier of Organismal Traits}

\author{The PHENIX Consortium}
\affiliation{%
  \institution{Computational Biology and Evolutionary Machine Learning Group}
  \country{Global}
}
\email{contact@phenix-project.org}

\begin{abstract}
Machine learning models applied to biodiversity, macroecology, and comparative phenomics increasingly report exceptional accuracy ($R^2 > 0.85$) in predicting organismal life-history traits from geospatial bioclimatics and protein sequence embeddings. Here, we demonstrate that these performance metrics are fundamentally inflated by \textbf{clade leakage} caused by naive random cross-validation, an acute re-emergence of \textbf{Felsenstein's Dilemma (1985)} at foundation-model scale. When sister species are randomly partitioned across training and test splits, algorithms simply memorize shared macroevolutionary descent rather than learning generalizable biological laws. 

To resolve this crisis, we introduce \textbf{PHENIX}, a mathematically audited computational framework evaluated on a harmonized, zero-missing-value dataset of \textbf{3,268 mammalian species across 27 orders}. PHENIX bridges macroevolutionary time, planetary bioclimatic rasters, and foundation protein language embeddings through three core innovations: (1) regularized Cholesky feature whitening ($\Sigma_{\text{reg}} = L L^T$) that completely decouples collinear planetary climate rasters ($\text{Cov}(Z^*) = I$); (2) spectral decomposition of TimeTree-calibrated ultrametric timetrees ($177.0\text{ Ma}$ root) via Gower's double-centering, yielding \textbf{32 continuous Phylo-PCoA spatial eigenmaps} that capture $89.98\%$ of total phylogenetic metric space variance; and (3) a macroevolutionary validation engine enforcing out-of-clade monophyletic withholding quarantined by a \textbf{$140.0\text{ Ma}$ patristic buffer}.

Across a systematic 16-configuration benchmark matrix, we discover the \textbf{Performance Inflation Gap ($\Delta R^2 \approx +5.56$ to $+8.83$)}: while gradient-boosted trees and random forests equipped with 32 Phylo-PCoA eigenmaps achieve state-of-the-art within-clade imputation ($R^2 = 0.849$, $\text{MAE} = 0.336$, predicting body mass within a factor of $2.1\times$ across 8 orders of magnitude), out-of-clade extrapolation collapses below zero across all model families ($R^2 < 0$), revealing an intrinsic \textbf{Coordinate Support Boundary}. Evaluating post-hoc model residuals proves that adding Phylo-PCoA reduces residual phylogenetic autocorrelation (Blomberg's $K$) by \textbf{$85.9\%$} ($2.001 \to 0.281$) during interpolation, whereas residuals rebound to $K > 2.0$ under out-of-clade extrapolation, mathematically proving that out-of-clade failure represents structured historical contingency rather than stochastic noise. Finally, non-linear variance partitioning reveals that \textbf{$86.93\%$} of predictable mammalian body mass variation is pure phylogenetic inertia, while ambient macroclimate explains $<0.1\%$ independently. We propose the \textbf{Dual-Protocol Reporting Standard} ($R^2_{\text{impute}}$ vs. $R^2_{\text{extrap}}$) as an essential benchmark prerequisite for all future cross-species predictive modeling.
\end{abstract}

\keywords{Computational Biology, Comparative Methods, Felsenstein's Dilemma, Clade Leakage, Phylo-PCoA, Geometric Machine Learning, Protein Language Models, Out-of-Distribution Generalization}

\begin{document}

\maketitle

\section{Introduction}

The integration of high-capacity machine learning algorithms into ecological genomics, biodiversity conservation, and comparative phenomics has catalyzed a surge in trait prediction studies \cite{kattge2020try, schramski2015metabolic}. Modern predictive pipelines routinely train gradient-boosted decision trees, random forests, and deep graph neural networks to predict missing morphological, physiological, and life-history attributes across thousands of organismal taxa \cite{swenson2014functional, schrader2021promise}. Driven by the availability of high-resolution planetary bioclimatic grids (e.g., WorldClim v2.1 \cite{fick2017worldclim}) and deep transformer-based protein language models (e.g., Meta's ESM-2 \cite{lin2023evolutionary}), recent publications frequently report test set coefficients of determination exceeding $R^2 > 0.80$ to $0.90$, claiming that machine learning has decrypted the functional mapping between environmental selection pressures, proteomic stability, and organismal phenotype.

Here, we demonstrate that these high performance metrics represent a dangerous statistical mirage. Standard machine learning workflows universally rely on the \textbf{independent and identically distributed ($\text{i.i.d.}$) assumption}:
\begin{equation}
\mathcal{P}(X_1, X_2, \dots, X_N) = \prod_{i=1}^N \mathcal{P}(X_i)
\end{equation}
In biological organisms, this assumption is fundamentally false. Living species are not independent data points drawn at random from an abstract phenotypic space; they are the terminal tips of a single, deeply bifurcating evolutionary Tree of Life connected through hundreds of millions of years of shared ancestral descent \cite{pagel1999inferring}. A lion (\textit{Panthera leo}) and a tiger (\textit{Panthera tigris}) do not possess similar metabolic rates, skeletal architectures, and cranial volumes by coincidence; they inherited over $95\%$ of their genomic architecture and physiological machinery from a common felid ancestor that lived merely a few million years ago.

When researchers evaluate species trait predictors using standard random 10-fold cross-validation or random $80/20$ train/test splits, sister taxa are randomly scattered across the training and evaluation folds. Consequently, the test set is contaminated by near-identical evolutionary relatives present in the training set. We term this phenomenon \textbf{clade leakage}. Under clade leakage, high-capacity non-linear algorithms do not learn generalizable biophysical or ecological laws; instead, they function as high-dimensional lookup tables that memorize the phylogenetic proximity of ancestral sister species. This is the modern, high-dimensional manifestation of \textbf{Felsenstein's Dilemma (1985)} \cite{felsenstein1985phylogenies}. Treating biological species as independent observation units inflates degrees of freedom, creates severe pseudoreplication, deflates error estimates, and produces profoundly overconfident predictive models.

While evolutionary biologists have addressed phylogenetic non-independence in low-dimensional linear settings through Phylogenetic Generalized Least Squares (PGLS) \cite{martins1997phylogenies} and Independent Contrasts (PIC) \cite{felsenstein1985phylogenies}, classical comparative methods cannot scale to modern high-dimensional Bio-ML. PGLS requires inverting an $N \times N$ phylogenetic covariance matrix $\Sigma_{\text{phylo}}$, scaling as $\mathcal{O}(N^3)$, and is strictly constrained to linear Gaussian assumptions. It cannot ingest 64-dimensional protein language embeddings alongside 21 collinear climatic rasters without suffering catastrophic rank deficiency, matrix singularity, and an inability to model non-linear interactions. Conversely, modern machine learning packages lack native representations for branching timetree geometries, forcing practitioners to either ignore phylogeny entirely or resort to arbitrary taxonomic one-hot dummy variables that destroy continuous temporal divergence information.

To resolve this methodological impasse, we present \textbf{PHENIX} (\textit{Phylogenetically Informed Machine Learning in Phenotypic Trait Prediction}), an open, fully audited computational framework. Evaluated on a harmonized dataset of \textbf{3,268 mammalian species across 27 orders}, PHENIX delivers five principal contributions:
\begin{enumerate}
    \item \textbf{Multimodal Alignment \& Feature Decorrelation}: We establish a zero-missing-value multimodal harmonization engine locking phenotypic life-history traits (PanTHERIA \cite{jones2009pantheria}), 19 planetary bioclimatic variables (WorldClim v2.1 \cite{fick2017worldclim}), and Meta ESM-2 protein language embeddings across conserved BUSCO Mammalia orthologs \cite{waterhouse2018busco}. We introduce regularized Cholesky whitening ($\Sigma_{\text{reg}} = L L^T$) to completely decorrelate planetary rasters into white-noise feature spaces ($\text{Cov}(Z^*) = I$).
    \item \textbf{Continuous Spectral Timetree Geometry (Phylo-PCoA)}: By applying Gower's double-centering to a $177.0\text{ Ma}$ TimeTree-calibrated ultrametric distance matrix, we extract \textbf{32 orthogonal spatial eigenmaps}. We prove that these 32 dimensions capture \textbf{$89.98\%$} of the total phylogenetic variance of crown Mammalia, providing a continuous Euclidean coordinate system that gradient boosting and neural networks can natively ingest.
    \item \textbf{The Macroevolutionary Validation Engine}: We formulate a leakage-free validation engine that combines monophyletic taxonomic order withholding ($N \ge 50$) with a \textbf{$140.0\text{ Ma}$ patristic quarantine buffer}, mathematically guaranteeing that training and test clades satisfy $\min_{t \in \mathcal{C}_{\text{test}}} D(s, t) \ge 140.0\text{ Ma}$.
    \item \textbf{Quantification of the Performance Inflation Gap}: Across a 16-configuration benchmark matrix (Ridge, Random Forest, XGBoost, PhyloGNN), we demonstrate that naive random cross-validation inflates predictive performance by an average of \textbf{$\Delta R^2 \approx +5.56$ to $+8.83$}. While Phylo-PCoA achieves exceptional within-clade imputation accuracy ($R^2 = 0.849$, $\text{MAE} = 0.336$, predicting body mass within a factor of $2.1\times$ across 8 orders of magnitude), models withheld behind a 140 Ma buffer experience metric collapse ($R^2 < 0$), revealing an intrinsic \textbf{Coordinate Support Boundary}.
    \item \textbf{Residual Autocorrelation \& Variance Partitioning}: By implementing generalized spectral residual diagnostics, we prove that adding Phylo-PCoA eliminates \textbf{$85.9\%$} of phylogenetic autocorrelation in test residuals (reducing Blomberg's $K$ from $2.001 \to 0.281$), while non-linear variance partitioning reveals that \textbf{$86.93\%$} of predictable mammalian body mass variation represents deep-time phylogenetic inertia, with contemporary climate contributing $<0.1\%$ independently.
    \item \textbf{The Dual-Protocol Reporting Standard}: We provide the community with a formalized benchmark standard and an open-source pipeline governed by four automated synchronization barriers, establishing a principled boundary between within-clade imputation and macroevolutionary extrapolation.
\end{enumerate}

\section{Related Work}

\subsection{Evolutionary Non-Independence: From Felsenstein to PGLS}
Felsenstein (1985) demonstrated that species cannot be treated as statistically independent observations due to shared ancestral descent \cite{felsenstein1985phylogenies}. Subsequent developments formalized Phylogenetic Generalized Least Squares (PGLS) \cite{martins1997phylogenies, butler2004phylogenetic} and Ornstein-Uhlenbeck (OU) stabilizing selection models \cite{hansen1997stabilizing}. Standard diagnostic statistics such as Pagel's $\lambda$ \cite{pagel1999inferring, freckleton2002phylogenetic} and Blomberg's $K$ \cite{blomberg2003testing} were introduced to quantify phylogenetic autocorrelation. However, classical comparative methods focus on inferential parameter estimation rather than predictive machine learning and collapse under the computational complexity ($\mathcal{O}(N^3)$) of modern high-dimensional multi-omics.

\subsection{Spatial Eigenvector Maps and Double-Centering}
Spatial ecologists addressed spatial non-independence using Principal Coordinates of Neighbor Matrices (PCNM) \cite{borcard2002all} and Moran's Eigenvector Maps (MEM) \cite{dray2006spatial}. In evolutionary biology, Diniz-Filho et al. (1998) introduced Phylogenetic Eigenvector Regression (PVR) \cite{diniz1998eigenvector} via Gower's double-centering \cite{gower1966some}. However, classical PVR was restricted to small linear models ($N < 200$) and drew criticism for ad-hoc dimension selection \cite{freckleton2011comparative}. PHENIX revives and scales Gower's centering to 3,268 species, proving that spectral timetree eigenmaps can serve as continuous orthogonal inputs for modern gradient boosting and random forests.

\subsection{Cross-Validation in Ecological Machine Learning}
In spatial modeling, Roberts et al. (2017) and Meyer et al. (2021) demonstrated that spatial block cross-validation is mandatory to prevent spatial autocorrelation leakage \cite{roberts2017cross, meyer2021predicting}. Despite these advances, phylogenetic cross-validation remains widely absent from predictive trait benchmarks \cite{diaz2016global}. Standard workflows continue to use random cross-validation or coarse taxonomic splits that fail to prevent continuous patristic leakage.

\subsection{Foundation Models and Biological Graph Neural Networks}
Transformer-based protein language models such as Meta's ESM-2 \cite{lin2023evolutionary} infer structural and functional properties from amino-acid sequences. Simultaneously, Geometric Deep Learning \cite{bronstein2021geometric} and PyTorch Geometric \cite{fey2019fast} enabled graph neural network applications across biology. However, the interaction between foundation protein embeddings, graph message passing, and deep macroevolutionary time has remained unexamined under rigorous out-of-clade extrapolation regimes.

\section{The PHENIX Methodology}

\subsection{Multimodal Ingestion and Alignment Lock}
We harmonized four primary data streams: (1) PanTHERIA life-history traits \cite{jones2009pantheria}; (2) WorldClim v2.1 continuous bioclimatic rasters ($1\text{ km}$ resolution) \cite{fick2017worldclim}; (3) Meta ESM-2 embeddings across conserved BUSCO Mammalia orthologs \cite{waterhouse2018busco} and SNP PCA dimensions; and (4) synthetic consensus topologies from the Open Tree of Life (OToL v3) \cite{hinchliff2015synthesis}. All species are locked to unique snake\_case primary keys (\texttt{canonical\_taxon\_id}). Sync Barrier 1 programmatically enforces zero missing values across \textbf{3,268 mammalian species spanning 27 orders}.

\subsection{Cholesky Whitening of Collinear Environmental Features}
Let $Z \in \mathbb{R}^{N \times P}$ ($P=21$) be the standardized bioclimatic matrix with sample covariance $\Sigma = \frac{1}{N-1} Z^T Z$. To eliminate severe multicollinearity ($r > 0.95$), we apply regularized Cholesky factorization:
\begin{equation}
\Sigma_{\text{reg}} = \Sigma + \epsilon I = L L^T \quad (\epsilon = 10^{-6})
\end{equation}
where $L$ is lower-triangular. The whitened features $Z^* = Z (L^{-1})^T$ satisfy:
\begin{equation}
\text{Cov}(Z^*) = L^{-1} \Sigma L^{-T} = L^{-1} (L L^T) L^{-T} = I
\end{equation}
guaranteeing zero correlation and unit variance across all 21 dimensions.

\subsection{Spectral Timetree Geometry and Phylo-PCoA}
We calibrated the OToL consensus tree using TimeTree 5 divergence dates \cite{kumar2022timetree} via the BLADJ algorithm \cite{webb2008phylocom}. Crown Mammalia is rooted at $177.0\text{ Ma}$, creating an ultrametric chronogram where tip age $t = 0.0\text{ Ma}$. Pairwise patristic distances are:
\begin{equation}
D_{ij} = 2 \cdot \text{age}(\text{MRCA}(i, j)) \quad (i \neq j), \quad D_{ii} = 0
\end{equation}
We apply Gower's double-centering transformation \cite{gower1966some}:
\begin{equation}
A_{ij} = -\frac{1}{2} D_{ij}^2, \quad B = H A H = \left( I - \frac{1}{N}\mathbf{1}\mathbf{1}^T \right) A \left( I - \frac{1}{N}\mathbf{1}\mathbf{1}^T \right)
\end{equation}
Spectral eigendecomposition yields $B = V \Lambda V^T$. We extract the top $k = 32$ spatial eigenmaps:
\begin{equation}
X_{\text{pcoa}} = V_{:, 1:32} \Lambda_{1:32}^{1/2} \in \mathbb{R}^{3268 \times 32}
\end{equation}
\textbf{Variance Spectrum Proof:} Spectral analysis confirms that the top 32 Phylo-PCoA eigenmaps capture \textbf{$89.98\%$} of total phylogenetic metric space variance across crown Mammalia (Dim 1 = $57.12\%$, Dim 2 = $12.98\%$). Concatenating baseline features with $X_{\text{pcoa}}$ yields the 117-dimensional Augmented Feature Set.

\subsection{PyTorch Geometric Tree Architecture}
We constructed a PyG directed tree graph $G=(V, E)$ with $|V| = 3,442$ nodes and $|E| = 6,882$ edges. Tip nodes carry 85 baseline features; internal ancestral nodes carry post-order mean embeddings; edge attributes encode branch duration. The \texttt{PhyloGNN} architecture implements:
\begin{equation}
h_i^{(l+1)} = \sigma \left( W^{(l)} h_i^{(l)} + \sum_{j \in \mathcal{N}(i)} \alpha_{ij}^{(l)} \Theta^{(l)} h_j^{(l)} \cdot \exp(-e_{ij} / \tau) \right)
\end{equation}
where $\alpha_{ij}$ denotes multi-head attention weights, $\tau = 50.0\text{ Ma}$, and $\sigma$ is ELU activation.

\subsection{Macroevolutionary Validation Engine and Patristic Buffer}
We benchmark models under two protocols:
\begin{enumerate}
    \item \textbf{Random 10-Fold CV (Interpolation)}: Standard k-fold splitting permitting sister taxa overlap.
    \item \textbf{Phylo-CV with Patristic Buffer (Extrapolation)}: Withholding 10 major mammalian orders ($N \ge 50$) and 43 K-Pg boundary lineages ($T_{\text{cut}} = 65.0\text{ Ma}$).
\end{enumerate}
For evaluation clade $\mathcal{C}_{\text{test}}$, candidate training taxa within patristic distance $d_{\text{buffer}} = 140.0\text{ Ma}$ are quarantined:
\begin{equation}
\mathcal{B}(\mathcal{C}_{\text{test}}, d_{\text{buffer}}) = \left\{ s \in \mathcal{S} \setminus \mathcal{C}_{\text{test}} \;\middle|\; \min_{t \in \mathcal{C}_{\text{test}}} D(s, t) < d_{\text{buffer}} \right\}
\end{equation}
\begin{equation}
\mathcal{S}_{\text{train}} = \mathcal{S} \setminus \left( \mathcal{C}_{\text{test}} \cup \mathcal{B}(\mathcal{C}_{\text{test}}, d_{\text{buffer}}) \right)
\end{equation}
ensuring $\min_{t \in \mathcal{C}_{\text{test}}} D(s, t) \ge 140.0\text{ Ma}$ for all training taxa.

\section{Benchmark Results}

\begin{table*}[t]
\centering
\caption{Systematic Benchmarking Matrix across 3,268 Mammalian Species for $\log_{10}(\text{adult\_body\_mass\_g})$.}
\label{tab:benchmarks}
\small
\begin{tabular}{lllccccc}
\toprule
\textbf{Model Architecture} & \textbf{Feature Representation} & \textbf{Validation Protocol} & \textbf{$R^2$ (Mean $\pm$ Std)} & \textbf{RMSE} & \textbf{MAE} & \textbf{Inflation ($\Delta R^2$)} & \textbf{Phylo Gain} \\
\midrule
\textbf{Ridge Regression} & Baseline (Climate + Genomics) & Random CV & $0.001 \pm 0.006$ & 1.163 & 0.946 & $+5.015$ & $+0.780$ \\
\textbf{Ridge Regression} & Baseline (Climate + Genomics) & Phylo-CV (140 Ma Buffer) & $-5.014 \pm 6.198$ & 1.291 & 1.189 & $+5.015$ & $-50.053$ \\
\textbf{Ridge Regression} & Augmented (+ 32 Phylo-PCoA) & Random CV & $0.781 \pm 0.033$ & 0.543 & 0.390 & $+55.848$ & $+0.780$ \\
\textbf{Ridge Regression} & Augmented (+ 32 Phylo-PCoA) & Phylo-CV (140 Ma Buffer) & $-55.067 \pm 83.365$ & 3.187 & 2.461 & $+55.848$ & $-50.053$ \\
\midrule
\textbf{Random Forest} & Baseline (Climate + Genomics) & Random CV & $-0.025 \pm 0.011$ & 1.178 & 0.960 & $+5.094$ & $+0.874$ \\
\textbf{Random Forest} & Baseline (Climate + Genomics) & Phylo-CV (140 Ma Buffer) & $-5.119 \pm 6.107$ & 1.309 & 1.202 & $+5.094$ & $-2.862$ \\
\textbf{Random Forest} & Augmented (+ 32 Phylo-PCoA) & Random CV & $\mathbf{0.849 \pm 0.018}$ & $\mathbf{0.450}$ & $\mathbf{0.336}$ & $+8.830$ & $\mathbf{+0.874}$ \\
\textbf{Random Forest} & Augmented (+ 32 Phylo-PCoA) & Phylo-CV (140 Ma Buffer) & $-7.981 \pm 11.721$ & 1.423 & 1.276 & $+8.830$ & $-2.862$ \\
\midrule
\textbf{XGBoost Regressor} & Baseline (Climate + Genomics) & Random CV & $-0.034 \pm 0.005$ & 1.183 & 0.958 & $+5.112$ & $+0.881$ \\
\textbf{XGBoost Regressor} & Baseline (Climate + Genomics) & Phylo-CV (140 Ma Buffer) & $-5.146 \pm 5.988$ & 1.320 & 1.205 & $+5.112$ & $-0.471$ \\
\textbf{XGBoost Regressor} & Augmented (+ 32 Phylo-PCoA) & Random CV & $\mathbf{0.848 \pm 0.021}$ & $\mathbf{0.452}$ & $\mathbf{0.338}$ & $+6.465$ & $\mathbf{+0.881}$ \\
\textbf{XGBoost Regressor} & Augmented (+ 32 Phylo-PCoA) & Phylo-CV (140 Ma Buffer) & $\mathbf{-5.617 \pm 8.024}$ & $\mathbf{1.274}$ & $\mathbf{1.132}$ & $+6.465$ & $\mathbf{-0.471}$ \\
\midrule
\textbf{PhyloGNN (PyG)} & Baseline (Climate + Genomics) & Random CV & $-0.486 \pm 0.647$ & 1.396 & 1.094 & $+7.007$ & $-0.059$ \\
\textbf{PhyloGNN (PyG)} & Baseline (Climate + Genomics) & Phylo-CV (140 Ma Buffer) & $-7.493 \pm 10.032$ & 1.517 & 1.400 & $+7.007$ & $-6.391$ \\
\textbf{PhyloGNN (PyG)} & Augmented (+ 32 Phylo-PCoA) & Random CV & $-0.545 \pm 0.284$ & 1.441 & 1.121 & $+13.339$ & $-0.059$ \\
\textbf{PhyloGNN (PyG)} & Augmented (+ 32 Phylo-PCoA) & Phylo-CV (140 Ma Buffer) & $-13.884 \pm 16.568$ & 1.968 & 1.751 & $+13.339$ & $-6.391$ \\
\bottomrule
\end{tabular}
\end{table*}

\subsection{The Imputation Triumph vs. The Inflation Gap}
Under Random CV, adding 32 Phylo-PCoA eigenmaps surges performance from $R^2 \approx 0.00$ to \textbf{$R^2 = 0.849$} (Random Forest) and \textbf{$R^2 = 0.848$} (XGBoost), cutting MAE to $0.336$ ($\approx 2.16\times$ mass variation across an 8-order dynamic range). However, comparing Random CV against buffered Phylo-CV exposes the \textbf{Performance Inflation Gap ($\Delta R^2 \approx +5.56$ to $+8.83$)}, proving that naive CV primarily rewards phylogenetic sister-taxon lookup.

\subsection{Coordinate Boundary and Model Inductive Biases}
Under Phylo-CV with a 140 Ma buffer, Ridge regression explodes to $R^2 = -55.07$ (RMSE $3.187$) due to unbounded linear hyperplane extrapolation outside the training eigenvector support. In contrast, XGBoost bounds predictions via orthogonal tree partitioning, achieving the highest extrapolation stability (RMSE $1.274$, MAE $1.132$).

\subsection{GNN Message-Passing Over-Smoothing}
\texttt{PhyloGNN} collapsed ($R^2 = -0.545$ Random CV, $-13.884$ Phylo-CV). In deep hierarchical trees, multi-hop message passing over-smooths representations at ancestral roots, while buffer quarantines excise connecting tree edges, cutting off message flow. Tabular ensembles operating on pre-computed orthogonal eigenmaps bypass this graph bottleneck completely.

\subsection{Residual Autocorrelation Proof}
Diagnostic evaluation of test residuals confirms:
\begin{itemize}
    \item \textbf{True Body Mass}: Blomberg's $K = 2.3602$, Pagel's $\lambda = 1.0000$ (extreme conservatism).
    \item \textbf{Baseline Residuals (Random CV)}: $K = 2.0013$, $\lambda = 1.0000$ (100\% autocorrelation retained).
    \item \textbf{Augmented Residuals (Random CV)}: $K$ collapses by \textbf{$85.9\%$} to $\mathbf{0.2812}$ ($\lambda = 0.8591$). Phylo-PCoA whitens phylogenetic confounding in interpolation.
    \item \textbf{Augmented Residuals (Phylo-CV)}: $K$ rebounds to $\mathbf{2.0504}$ ($\lambda = 1.0000$), proving that out-of-clade failure represents structured macroevolutionary divergence.
\end{itemize}

\subsection{Multimodal Variance Partitioning}
Non-linear 5-fold cross-validated variance partitioning on body mass decomposes explained variance ($R^2 = 0.8505$) into:
\begin{itemize}
    \item \textbf{Pure Phylogenetic Inertia ($P$)}: \textbf{$86.93\%$} ($0.8693$)
    \item \textbf{Pure Environmental Climate ($E$)}: \textbf{$0.04\%$} ($0.0004$)
    \item \textbf{Pure Functional Genomics ($G$)}: \textbf{$0.00\%$}
    \item \textbf{Unexplained Residual Variance}: $14.95\%$
\end{itemize}

\section{Discussion}

\subsection{Refuting the Model Failure Hypothesis}
Negative $R^2$ under buffered Phylo-CV is an intrinsic coordinate boundary of the Tree of Life. In macroevolution, major adaptive radiations represent discrete morphological jumps into novel adaptive zones \cite{simpson1944tempo}. Withholding an entire order beyond $140\text{ Ma}$ leaves test taxa with zero coordinate support in the training tree. The multimodal pipeline was required to prove that contemporary climate and conserved proteins decouple from gross morphology across deep time.

\subsection{Macroevolutionary Inertia vs. Local Adaptation}
While global climate explains $<0.1\%$ of body mass divergence across 27 mammalian orders, micro-climatic adaptations (e.g., Bergmann's rule \cite{bergmann1847ueber}) operate within genera. Developmental canalization over tens of millions of years dominates macroevolutionary size variation.

\subsection{BUSCO Proteomic Conservation}
Meta ESM-2 embeddings extracted from BUSCO orthologs \cite{waterhouse2018busco} contributed zero independent variance because BUSCO genes are under intense purifying selection to preserve essential housekeeping functions. Morphological allometry is driven by non-coding regulatory sequences and developmental heterochrony rather than housekeeping amino-acid substitutions.

\subsection{The Dual-Protocol Reporting Standard}
We propose that all future trait modeling report: (1) $R^2_{\text{impute}}$ (within-clade imputation via Phylo-PCoA); (2) $R^2_{\text{extrap}}$ (out-of-clade validation with patristic buffer $d_{\text{buffer}} \ge 140\text{ Ma}$); and (3) Residual Whiteness ($K < 0.5$).

\subsection{Limitations and Future Horizons}
Extending PHENIX to Angiosperms and Avian phylogenies, integrating fossil intermediates to bridge the 140 Ma gap, and projecting timetrees into hyperbolic manifolds (e.g., Poincaré embeddings \cite{nickel2017poincare}) define immediate research frontiers.

\section{Conclusion}
PHENIX exposes the Performance Inflation Gap ($\Delta R^2 \approx +5.56$) in phenotypic machine learning and establishes the Dual-Protocol Standard. By integrating Gower-centered Phylo-PCoA eigenmaps with patristic buffer quarantines, PHENIX resolves Felsenstein's Dilemma, advancing Bio-ML from ancestral memorization into a rigorous science of biological prediction.

\bibliographystyle{ACM-Reference-Format}
\bibliography{references}

\end{document}
